\chapter{Methodology}
\label{ch:methodology}

This chapter details the methodology used to develop, simulate, and analyze the decentralized edge machine learning system for predicting environmental path loss in indoor LoRaWAN networks. Aligning with the constraints of TinyML and low-power wide-area networks, the proposed framework leverages a highly optimized neural network architecture, naturally non-IID spatial data partitioning, and advanced federated optimization algorithms to train models collaboratively without transmitting raw measurement data.

% ─────────────────────────────────────────────────────────────────────────────
% 4.1 RESEARCH SETUP AND MULTI-LAYER ARCHITECTURE
% ─────────────────────────────────────────────────────────────────────────────
\section{Research Setup and Multi-Layer Architecture}
\label{sec:research-setup}

\subsection{Overview of the Multi-Layer Architecture}
The proposed system is structured as a three-layer architecture designed to support decentralized, collaborative model training on battery-powered edge nodes:
\begin{enumerate}
    \item \textbf{Offline Model Preparation (Layer 1):} Prior to deployment, raw environmental telemetry and radio logs from historical campaigns are preprocessed. A baseline Neural Network model is defined, and feature scaling parameters (means and standard deviations) are computed.
    \item \textbf{Local Edge Intelligence (Layer 2):} The edge nodes (deployed microcontrollers) perform continuous sensor data acquisition, feature normalization, local path loss inference, and event-driven transmission scheduling. Crucially, when triggered, the nodes perform local training on a small circular buffer of recent measurements to compute weight updates (deltas) which are then quantized and transmitted.
    \item \textbf{Federated Server Orchestration (Layer 3):} A central server receives the local updates via gateway webhooks, aggregates them using federated optimization algorithms (e.g., FedAvg, FedProx, SCAFFOLD, FedYogi, or FedDyn), updates the global weights, and downlinks the new global model back to the edge nodes.
\end{enumerate}

\subsection{Hardware Platform and Sensor Topology}
The hardware platform consists of six physical end devices (ED0–ED5) deployed across the 8th floor of the Hölderlinstraße Campus, University of Siegen. The key specifications of the hardware prototype are:
\begin{itemize}
    \item \textbf{Microcontroller (Arduino MKR WAN 1310):} Features a Microchip SAMD21 32-bit ARM Cortex-M0+ processor running at 48 MHz, with 32 KB of SRAM and 256 KB of Flash memory.
    \item \textbf{LoRa Module (Murata CMWX1ZZABZ-078):} Embeds a Semtech SX1276 LoRa transceiver. The devices are configured for the EU868 band, Class A operation, using Over-the-Air Activation (OTAA).
    \item \textbf{Sensor Suite:} Each node is integrated with three digital environmental sensors interfacing via I\textsuperscript{2}C:
    \begin{itemize}
        \item \textbf{Sensirion SCD4x:} Captures carbon dioxide ($\text{CO}_2$) concentration (range: 400--5000 ppm), ambient temperature (range: $-10$ to $+60\,^{\circ}\text{C}$), and relative humidity (range: 0--100 \%RH).
        \item \textbf{Bosch BME280:} Monitors barometric pressure (range: 300--1100 hPa).
        \item \textbf{Sensirion SPS30:} Measures fine particulate matter ($\text{PM}_{2.5}$) concentration (range: 0--1000 $\mu\text{g/m}^3$).
    \end{itemize}
\end{itemize}

% ─────────────────────────────────────────────────────────────────────────────
% 4.2 ENVIRONMENTAL DATA PIPELINE AND SPATIAL PARTITIONING
% ─────────────────────────────────────────────────────────────────────────────
\section{Environmental Data Pipeline and Spatial Partitioning}
\label{sec:data-pipeline}

\subsection{The Environmental Telemetry Dataset}
The dataset used in this work originates from the indoor LoRaWAN measurement campaign conducted by Obiri and Van Laerhoven~\cite{obiri2024comprehensive}. The campaign recorded environmental observations alongside LoRaWAN radio telemetry across six physical end devices (ED0–ED5). The raw data logged in the InfluxDB database contains 2,079,534 telemetry records.

\subsection{Data Cleaning and Quality Rectification}
To prepare the dataset for model training, a deterministic preprocessing pipeline was executed:
\begin{enumerate}
    \item \textbf{Telemetry Alignment:} Atmospheric pressure values recorded by the BME280 sensor were multiplied by 3.125 to restore them to the physical hPa scale.
    \item \textbf{Null Handling:} Rows containing null values in the signal-to-noise ratio (SNR) or packet frame count columns were dropped.
    \item \textbf{Outlier Filtering:} The target variable, expected path loss ($\text{exp\_pl}$ in dB), is defined as:
    \begin{equation{
        \text{exp\_pl} = P_{\text{TX}} - P_{\text{RX}} + G_{\text{TX}} + G_{\text{RX}}
    \end{equation}
    where $P_{\text{TX}} = 14$ dBm is the transmission power, $P_{\text{RX}}$ is the RSSI, and $G_{\text{TX}}, G_{\text{RX}}$ are antenna gains. A strict range filter was applied to restrict $\text{exp\_pl}$ to a physically plausible range of $[50, 200]$ dB.
\end{enumerate}
This pipeline yielded a final usable dataset of \textbf{2,079,528} records. The preprocessing stages and row counts are summarized in Table~\ref{tab:data-cleaning}.

\begin{table}[H]
  \centering
  \begin{tabular}{lr}
    \toprule
    \textbf{Preprocessing Stage} & \textbf{Remaining Rows} \\
    \midrule
    Raw Telemetry Logs (TTN Database)           & 2,079,534 \\
    Pressure Scaling Correction ($\times 3.125$) & 2,079,534 \\
    Null SNR \& Frame Count Drop                & 2,079,534 \\
    Path Loss Outlier Filter ($[50, 200]$\,dB)   & 2,079,528 \\
    \textbf{Final Preprocessed Dataset}        & \textbf{2,079,528} \\
    \bottomrule
  \end{tabular}
  \caption{Data preprocessing pipeline and row counts at each stage.}
  \label{tab:data-cleaning}
\end{table}

\subsection{Input Feature Selection and Target Definition}
The regression model utilizes nine input features ($D=9$) to predict the expected path loss ($\text{exp\_pl}$):
\begin{itemize}
    \item \textbf{Spatial and Geometric Features (3):}
    \begin{itemize}
        \item $\text{log\_distance}$ ($10\log_{10}(d)$): Logarithmic physical distance $d$ between the end device and the gateway.
        \item $W_{\text{brick}}$: The count of physical brick walls intersecting the direct line-of-sight path.
        \item $W_{\text{wood}}$: The count of physical wooden partitions intersecting the line-of-sight path.
    \end{itemize}
    \item \textbf{Environmental Features (5):}
    \begin{itemize}
        \item $\text{co}_2$: Carbon dioxide concentration (ppm).
        \item $\text{humidity}$: Relative humidity (\%RH).
        \item $\text{pm}_{2.5}$: Particulate matter concentration ($\mu\text{g/m}^3$).
        \item $\text{pressure}$: Atmospheric pressure (hPa).
        \item $\text{temperature}$: Ambient temperature ($^{\circ}$C).
    \end{itemize}
    \item \textbf{Radio and Link Quality Features (1):}
    \begin{itemize}
        \item $\text{snr}$: The signal-to-noise ratio (dB) of the received packet.
    \end{itemize}
\end{itemize}

\subsection{Naturally Non-IID Spatial Partitioning}
To simulate real-world conditions, the preprocessed dataset was partitioned using a natural \textit{non-IID by location} split. Rather than allocating samples randomly, the data was grouped by the originating physical device's identifier ($\text{device\_id} \in \{\text{ED0}, \dots, \text{ED5}\}$). This creates six distinct local datasets, each representing the unique environmental and spatial characteristics of its respective room. Table~\ref{tab:client-data} lists the sample counts assigned to each virtual client.

\begin{table}[H]
  \centering
  \begin{tabular}{lllr}
    \toprule
    \textbf{Client Index} & \textbf{Device ID} & \textbf{Room Placement} & \textbf{Local Samples} \\
    \midrule
    Client 1 & ED0 & Office A (Near Window) & 277,531 \\
    Client 2 & ED1 & Main Corridor           & 274,660 \\
    Client 3 & ED2 & Kitchen                 & 278,728 \\
    Client 4 & ED3 & Office B (Interior)     & 272,897 \\
    Client 5 & ED4 & Stairwell (North-End)   & 275,528 \\
    Client 6 & ED5 & Meeting Room            & 284,278 \\
    \midrule
    \textbf{Total} & & & \textbf{2,079,528} \\
    \bottomrule
  \end{tabular}
  \caption{Sample distribution across virtual clients (Non-IID partitioning by location).}
  \label{tab:client-data}
\end{table}

% ─────────────────────────────────────────────────────────────────────────────
% 4.3 ON-DEVICE TINYML MODEL ARCHITECTURE
% ─────────────────────────────────────────────────────────────────────────────
\section{On-Device TinyML Model Architecture}
\label{sec:tinyml-arch}

\subsection{Neural Network Topology}
To satisfy the memory constraints of the SAMD21 microcontroller and minimize transmission overhead over LoRaWAN, a shallow Feedforward Neural Network (FNN) was designed. The architecture, denoted as Dense($9{\to}8{\to}1$), consists of:
\begin{enumerate}
    \item \textbf{Input Layer ($D=9$):} Accepts the standardized 9-dimensional spatial, environmental, and link quality features.
    \item \textbf{Hidden Layer ($H=8$):} Eight fully connected neurons using the Rectified Linear Unit (ReLU) activation function:
    \begin{equation}
        f(x) = \max(0, x)
    \end{equation}
    \item \textbf{Output Layer ($O=1$):} A single neuron with a linear activation function that outputs the predicted expected path loss ($\text{exp\_pl}$) in dB.
\end{enumerate}

\subsection{Design Justifications for Resource Constraints}
Selecting a minimal, shallow architecture is critical for several key reasons:
\begin{itemize}
    \item \textbf{LoRaWAN Payload Alignment:} An 89-parameter model represented in 8-bit quantized weights is only 89 bytes. This easily fits inside a single LoRaWAN transmission packet (which supports payloads up to 222 bytes), eliminating the need for multi-packet fragmentation and reducing network congestion.
    \item \textbf{TinyML Microcontroller Constraints:} The target Arduino hardware features only 32 KB of SRAM and 256 KB of Flash memory. The Dense($9{\to}8{\to}1$) model compiled into TensorFlow Lite format occupies less than 3 KB of memory, leaving ample space for the sensor acquisition drivers and LoRaWAN communication stack.
    \item \textbf{Prevention of Non-IID Overfitting:} Deeper networks are prone to memorizing local client noise in non-IID setups, which leads to poor generalization on other clients' data. A highly constrained architecture acts as a structural regularizer.
    \item \textbf{Universal Approximation Theorem:} According to the theorem, a single hidden layer with a non-linear activation function (such as ReLU) can approximate any continuous function on compact subsets of $\mathbb{R}^n$, provided it has sufficient width.
\end{itemize}

\subsection{Input Standardization and Model Export}
Neural networks require input scaling to ensure stable gradient updates. A \texttt{StandardScaler} is fitted on the training set, computing the mean $\mu_j$ and standard deviation $\sigma_j$ for each of the 9 features. The normalization parameters are exported and hardcoded in the Arduino sketch. Input features are standardized on-device prior to inference using:
\begin{equation}
    x^{\text{norm}}_j = \frac{x_j - \mu_j}{\sigma_j}
\end{equation}

\subsection{Mathematical Derivation of Trainable Parameters}
The total parameter count $N_{\text{params}}$ of a multi-layer perceptron with a single hidden layer is calculated as:
\begin{equation}
    N_{\text{params}} = (D \times H + H) + (H \times O + O)
\end{equation}
Substituting $D = 9$, $H = 8$, and $O = 1$:
\begin{equation}
    N_{\text{params}} = (9 \times 8 + 8) + (8 \times 1 + 1) = 80 + 9 = 89
\end{equation}
Thus, the model contains exactly 89 trainable weights and biases.

% ─────────────────────────────────────────────────────────────────────────────
% 4.4 FEDERATED OPTIMIZATION AND ALGORITHMIC DESIGN
% ─────────────────────────────────────────────────────────────────────────────
\section{Federated Optimization and Algorithmic Design}
\label{sec:fl-algorithms}

To address the challenges of spatial non-IID data distribution, five federated optimization algorithms were implemented and evaluated under identical conditions.

\subsection{Baseline Collaborative Learning (FedAvg)}
The baseline FedAvg algorithm aggregates local model weights by computing a weighted average based on client sample sizes:
\begin{equation}
    w^{t+1} = \sum_{k=1}^K \frac{n_k}{N} w_k^{t+1}
\end{equation}
where $n_k$ is the number of samples on client $k$, and $N = \sum_k n_k$. While effective on IID data, FedAvg suffers from client drift on non-IID data, dragging the global model towards local optima.

\subsection{Addressing Non-IID Local Drift (FedProx)}
FedProx addresses client drift by adding an isotropic $L_2$ regularization term to the local loss function. The local objective for client $k$ becomes:
\begin{equation}
    \mathcal{L}_k(w) = \mathcal{L}_{\text{MSE}}(w; \mathcal{D}_k) + \frac{\mu}{2} \|w - w^t\|^2
\end{equation}
where $w^t$ represents the global weights at the start of the round, and $\mu = 0.01$ is the proximal penalty coefficient. This term restricts the local updates from drifting too far from the global consensus.

\subsection{Inter-Client Gradient Variance Correction (SCAFFOLD)}
SCAFFOLD uses stateful control variates (server control $c$ and client control $c_k$) to estimate the client drift direction. The local gradient update step on the client is modified as:
\begin{equation}
    g_k(w) \leftarrow g_k(w) - c_k + c
\end{equation}
At the end of the round, the control variates are updated based on the weight changes:
\begin{align}
    c_k^{t+1} &= c_k^t - c + \frac{1}{E \eta} (w^t - w_k^{t+1}) \\
    c^{t+1} &= \frac{1}{K} \sum_{k=1}^K c_k^{t+1}
\end{align}
This correction vector offsets the local gradient step, aligning client trajectories back towards the global optimum.

\subsection{Server-Side Adaptive Learning (FedYogi)}
FedYogi is an adaptive server-side optimizer that replaces standard averaging with adaptive gradient updates. Instead of setting $w^{t+1} = w_{\text{avg}}$, the server computes a pseudo-gradient:
\begin{equation{
    \Delta w^t = w^t - \sum_{k=1}^K \frac{n_k}{N} w_k^{t+1}
\end{equation}
It then updates the server momentum $m_t$ and variance $v_t$ vectors:
\begin{align}
    m_t &= \beta_1 m_{t-1} + (1 - \beta_1) \Delta w^t \\
    v_t &= v_{t-1} + (1 - \beta_2) \text{sign}(\Delta w^t) (\Delta w^t)^2
\end{align}
The global model is updated as:
\begin{equation}
    w^{t+1} = w^t - \eta_s \frac{m_t}{\sqrt{v_t} + \epsilon}
\end{equation}
where $\beta_1 = 0.9$, $\beta_2 = 0.99$, $\eta_s = 0.01$, and $\epsilon = 10^{-3}$. FedYogi controls learning rates per parameter, preventing dominant clients from destabilizing convergence.

\subsection{Dynamic Client Regularization (FedDyn)}
FedDyn dynamically adjusts the local client objectives using a penalty term that adapts based on historic client updates. For client $k$, the local objective is:
\begin{equation}
    \mathcal{L}_k(w) = \mathcal{L}_{\text{MSE}}(w; \mathcal{D}_k) - \langle \nabla \mathcal{L}_k(w^t), w \rangle + \frac{\alpha}{2} \|w - w^t\|^2
\end{equation}
The gradient term is replaced by a running cumulative parameter drift vector $h_k$:
\begin{equation}
    \mathcal{L}_k(w) = \mathcal{L}_{\text{MSE}}(w; \mathcal{D}_k) + \alpha \langle w, -w^t + h_k^t \rangle + \frac{\alpha}{2} \|w - w^t\|^2
\end{equation}
where $\alpha = 0.01$. At the end of round $t$, the client updates its local state vector:
\begin{equation}
    h_k^{t+1} = h_k^t + (w_k^{t+1} - w^t)
\end{equation}
The server aggregates the local weights and applies the average state offset:
\begin{align}
    w_{\text{avg}}^{t+1} &= \sum_{k=1}^K \frac{n_k}{N} w_k^{t+1} \\
    w^{t+1} &= w_{\text{avg}}^{t+1} + \frac{1}{K}\sum_{k=1}^K h_k^{t+1}
\end{align}
By combining local constraints with memory of historical drift, FedDyn asymptotically aligns local training directions with the global objective.

% ─────────────────────────────────────────────────────────────────────────────
% 4.5 DEPLOYMENT FEASIBILITY AND CONSTRAINTS ANALYSIS
% ─────────────────────────────────────────────────────────────────────────────
\section{Deployment Feasibility and Constraints Analysis}
\label{sec:deployment-feasibility}

\subsection{Microcontroller SRAM and Flash Footprint}
To validate edge compatibility, the memory requirements of the Dense($9{\to}8{\to}1$) model were analyzed. The storage (Flash) and runtime (SRAM) footprints on the SAMD21 Cortex-M0+ are calculated as follows:
\begin{itemize}
    \item \textbf{Model Parameters (Flash):} The 89 weights and biases are stored as 32-bit floats, consuming:
    \begin{equation}
        89 \times 4\text{ bytes} = 356\text{ bytes}
    \end{equation}
    When converted to the TensorFlow Lite format, the serialized model schema introduces wrapper overhead, producing a \texttt{model.tflite} file of 2,704 bytes. This represents only 1.05\% of the SAMD21's 256 KB Flash limit.
    \item \textbf{Tensor Arena (SRAM):} The TensorFlow Lite Micro interpreter requires a contiguous memory block (the tensor arena) to allocate input, output, and intermediate layer activation tensors. Due to the small size of the activations ($9 \times 1$ input, $8 \times 1$ hidden, and $1 \times 1$ output), a tensor arena of 8,192 bytes (8 KB) is sufficient. This consumes 25% of the SAMD21's 32 KB SRAM, leaving 24 KB for the LoRaWAN stack and sensor code.
\end{itemize}

\subsection{LoRaWAN Payload and Quantized Transmission Scheme}
Transmission bandwidth is the primary bottleneck in LoRaWAN networks. To avoid fragmentation and minimize Time-on-Air (ToA), local model updates are quantized from 32-bit floats to 8-bit signed integers.
Given a parameter weight $w$, the quantized representation $q \in [-128, 127]$ is:
\begin{equation}
    q = \text{round}(w \times 127)
\end{equation}
This compression maps the 89 floats down to 89 bytes. In the uplink payload, weights are packed sequentially:
\begin{itemize}
    \item Hidden layer weights: $9 \times 8 = 72$ bytes.
    \item Hidden layer biases: 8 bytes.
    \item Output layer weights: 8 bytes.
    \item Output layer bias: 1 byte.
\end{itemize}
This results in a total model payload of exactly **89 bytes**. 

This payload design is compared with the benchmark established by Torres-Sanchez et al., who simulated an autoencoder model with 32 hidden neurons. Their model required transmitting 1,428 bytes per round. Under LoRaWAN restrictions, a 1.428 KB model cannot be transmitted in a single packet and must be fragmented into 7 to 28 separate frames depending on the Spreading Factor. In contrast, our 89-byte payload fits completely within a single packet across almost all Data Rates, reducing the network transmission overhead by a factor of 16.

\subsection{Time-on-Air and Link Budget Energy Modeling}
To quantify network overhead, the Time-on-Air (ToA) of the 89-byte model payload was modeled across different LoRaWAN Spreading Factors (SF) on the EU868 band (125 kHz bandwidth). The physical transmission time $T_{\text{packet}}$ is calculated as:
\begin{equation}
    T_{\text{packet}} = T_{\text{preamble}} + T_{\text{payload}}
\end{equation}
where $T_{\text{preamble}}$ is the preamble duration and $T_{\text{payload}}$ is the symbol count multiplied by the symbol duration.
Using the Semtech SX1276 link model, the ToA and associated energy consumption (assuming a typical transmission current of 25 mA at 3.3V) are detailed in Table~\ref{tab:toa-energy}.

\begin{table}[H]
  \centering
  \begin{tabular}{ccccc}
    \toprule
    \textbf{Spreading Factor} & \textbf{Data Rate} & \textbf{Max Payload} & \textbf{Time-on-Air (89\,B)} & \textbf{Energy (3.3V, 25\,mA)} \\
    \midrule
    SF7  & DR5 & 222 B & 143.6 ms  & 0.0118 J \\
    SF8  & DR4 & 222 B & 256.5 ms  & 0.0212 J \\
    SF9  & DR3 & 115 B & 451.1 ms  & 0.0372 J \\
    SF10 & DR2 & 51 B  & Fragmented (2 packets) & -- \\
    SF11 & DR1 & 51 B  & Fragmented (2 packets) & -- \\
    SF12 & DR0 & 51 B  & Fragmented (2 packets) & -- \\
    \bottomrule
  \end{tabular}
  \caption{LoRaWAN physical layer transmission metrics for the 89-byte model payload.}
  \label{tab:toa-energy}
\end{table}

As shown in Table~\ref{tab:toa-energy}, at SF7 through SF9, the model fits into a single uplink packet, consuming less than 0.04 Joules of energy. For SF10 to SF12, where the LoRaWAN maximum payload length drops to 51 bytes, the payload is split into two packets of 45 bytes and 44 bytes, which are sequentially. This analysis demonstrates that federated learning updates can be realistically supported within the constraints of public LoRaWAN networks.
